\documentclass[11pt,a4paper,UTF8]{ctexart}
\usepackage{geometry}
\geometry{left=18mm,right=18mm,top=24mm,bottom=24mm,headheight=22pt,footskip=12mm}
\usepackage{amsmath,amssymb,mathtools,bm,esint,extarrows}
\usepackage{xcolor}
\usepackage{tcolorbox}
\tcbuselibrary{skins,breakable}
\usepackage{fancyhdr}
\usepackage{titlesec}
\usepackage{hyperref}
\usepackage{tikz}
\usepackage{booktabs}
\usepackage{tabularx}

% --- Color Palette (Purple & DeepPink Academic Theme) ---
\definecolor{DeepPurple}{HTML}{4C1D95}
\definecolor{TitlePurple}{HTML}{581C87}
\definecolor{SubPurple}{HTML}{6B21A8}
\definecolor{DeepPink}{HTML}{FF1493}
\definecolor{LightPinkBg}{HTML}{FFF5F8}
\definecolor{LightPurpleBg}{HTML}{F8F5FF}
\definecolor{GoldAccent}{HTML}{D97706}
\definecolor{DarkText}{HTML}{1F2937}

\hypersetup{
    colorlinks=true,
    linkcolor=TitlePurple,
    citecolor=DeepPink,
    urlcolor=SubPurple,
    pdfborder={0 0 0}
}

% --- Typography & Spacing ---
\linespread{1.3}
\setlength{\parskip}{0.35em plus 0.1em minus 0.1em}
\setlength{\parindent}{0pt}

% --- Section Titles Styling ---
\titleformat{\section}
  {\Large\bfseries\color{TitlePurple}}{\thesection}{1em}{}
  [\vspace{1mm}{\color{TitlePurple!30}\hrule height 1pt}]
\titleformat{\subsection}
  {\large\bfseries\color{SubPurple}}{\thesubsection}{0.8em}{}
\titleformat{\subsubsection}
  {\normalsize\bfseries\color{DeepPurple}}{\thesubsubsection}{0.6em}{}

% --- tcolorbox Environments ---
% 1. Theorem / Formula / Method / Analysis Box (Deep Purple)
\newtcolorbox{thmbox}[1][]{
    enhanced,
    breakable,
    colback=LightPurpleBg,
    colframe=TitlePurple,
    arc=3pt,
    boxrule=0pt,
    leftrule=3.5pt,
    left=10pt,
    right=10pt,
    top=8pt,
    bottom=8pt,
    title=#1,
    coltitle=white,
    colbacktitle=TitlePurple,
    fonttitle=\bfseries\small,
    attach boxed title to top left={yshift=-2mm, xshift=4mm},
    boxed title style={boxrule=0pt, arc=2pt}
}

% 2. Solution / Formula / Derivation Box (DeepPink)
\newtcolorbox{solbox}[1][]{
    enhanced,
    breakable,
    colback=LightPinkBg,
    colframe=DeepPink,
    arc=3pt,
    boxrule=0pt,
    leftrule=3.5pt,
    left=10pt,
    right=10pt,
    top=8pt,
    bottom=8pt,
    title=#1,
    coltitle=white,
    colbacktitle=DeepPink,
    fonttitle=\bfseries\small,
    attach boxed title to top left={yshift=-2mm, xshift=4mm},
    boxed title style={boxrule=0pt, arc=2pt}
}

% 3. Learning Goals / Summary Box
\newtcolorbox{targetbox}[1][]{
    enhanced,
    breakable,
    colback=LightPurpleBg,
    colframe=DeepPurple,
    arc=3pt,
    boxrule=1pt,
    left=10pt,
    right=10pt,
    top=8pt,
    bottom=8pt,
    title=#1,
    coltitle=white,
    colbacktitle=DeepPurple,
    fonttitle=\bfseries\small,
    attach boxed title to top left={yshift=-2mm, xshift=4mm},
    boxed title style={boxrule=0pt, arc=2pt}
}

\begin{document}

% ------------------- 讲义封面 (Cover Page) -------------------
\begin{titlepage}
\thispagestyle{empty}
\begin{tikzpicture}[remember picture,overlay]
    \fill[DeepPurple] (current page.north west) rectangle ([yshift=-7cm]current page.north east);
    \draw[DeepPink, line width=3.5pt] ([yshift=-7cm]current page.north west) -- ([yshift=-7cm]current page.north east);
    \draw[GoldAccent, line width=1pt] ([yshift=-7.12cm]current page.north west) -- ([yshift=-7.12cm]current page.north east);
    
    \node[anchor=north east, opacity=0.12, text=white] at ([xshift=-1.5cm, yshift=-0.5cm]current page.north east) {
        \fontsize{70}{70}\selectfont\bfseries 2027
    };
\end{tikzpicture}

\vspace*{0.8cm}
\begin{center}
    {\color{white}\fontsize{26}{32}\selectfont\bfseries 2027 考研数学(二) 核心讲义}\\[0.6cm]
    {\color{white}\fontsize{13}{18}\selectfont\bfseries 全国硕士研究生招生考试 · 统考数学(二) 核心精讲全书}\\[1.5cm]
\end{center}

\vspace{3.5cm}
\begin{center}
    {\color{GoldAccent}\large\bfseries $\blacktriangleright$ 基础强化 · 核心题解 · 考点深水区 $\blacktriangleleft$}\\[0.8cm]
    {\color{TitlePurple}\fontsize{24}{28}\selectfont\bfseries 第十二章 一元积分学的应用(物理和经济)}\\[0.6cm]
    {\color{DeepPink}\large\bfseries —— 课后精选练习题与全过程推导 ——}\\[1.5cm]
\end{center}

\vfill

\begin{center}
\begin{tcolorbox}[
    colback=white,
    colframe=DeepPurple,
    width=0.88\textwidth,
    arc=4pt,
    boxrule=1.2pt,
    halign=center,
    drop shadow=gray!30
]
    \vspace{3mm}
    {\bfseries\large 编著团队：EscoffierZhou}\\[2.5mm]
    {\bfseries\color{TitlePurple} 目标院校：中国科学院大学 (UCAS) 计算机专硕 (22408)}\\[2.5mm]
    {\small\color{gray} 备考铁律：手算真题数字 · 错题白纸重做 · 408 客观题为王}\\[2.5mm]
    {\small 编制版本：2027 届首发版 · \today}
    \vspace{2mm}
\end{tcolorbox}
\end{center}
\vspace{0.8cm}
\end{titlepage}

% ------------------- 目录与页面主体配置 -------------------
\pagestyle{fancy}
\fancyhf{}
\fancyhead[L]{\small\color{gray}2027 考研数学(二) · 核心讲义系列}
\fancyhead[R]{\small\color{TitlePurple}\nouppercase{\leftmark}}
\fancyfoot[C]{\small\color{gray}— 第 \thepage\ 页 —}
\fancyfoot[R]{\small\color{gray}UCAS 22408}
\renewcommand{\headrulewidth}{0.5pt}

\pagenumbering{Roman}
\tableofcontents
\newpage
\pagenumbering{arabic}


\section{12.Homework}


\subsection{习题 12.1: 悬挂均匀铁链对折上提做功的双重解法}

一根长为 $L$、质量为 $M$ 的均质铁链垂直悬挂在天花板上。若将铁链的下端缓慢拉起，直到下端与悬挂的上端重合（即将铁链对折），求外力克服重力所做的功。

\textbf{\color{DeepPurple}【思路点拨】} 解法[1]采用积分微元法，以天花板悬点为坐标原点，向下为正建立坐标轴，分析下半段铁链中每个微元段被对折提升的高度，建立功微元并积分；解法[2]采用质心势能法，利用对折前后铁链各部分质心高度的变化，由外力做功等于重力势能增量直接求解


\begin{solbox}[分步推导与答题规范]
\textbf{[1]} 解法[1]（定积分微元分析法）：

取悬挂点为原点 $O$，铅直向下为 $x$ 轴正向。铁链全长分布在 $x \in [0, L]$。

均质铁链的线密度为 $\rho_0 = \frac{M}{L}$。

将铁链下端提起与上端重合的过程中，铁链的上半段 $[0, \frac{L}{2}]$ 保持原位不动，仅有下半段 $[\frac{L}{2}, L]$ 被翻折拉起。

对于下半段处于坐标 $x \in [\frac{L}{2}, L]$ 处的微元段 $\mathrm{d}x$，其质量为 $\mathrm{d}m = \rho_0\,\mathrm{d}x = \frac{M}{L}\,\mathrm{d}x$。

当对折完成时，原位于 $x$ 处的微元被折叠至上半段关于中点 $\frac{L}{2}$ 对称的位置，即新坐标为 $x' = L - x$。

因此，该微元被提升的垂直高度为：


\begin{equation*}
\Delta h(x) = x - x' = x - (L - x) = 2x - L
\end{equation*}


外力克服重力提升该微元所做的功为：


\begin{equation*}
\mathrm{d}W = \mathrm{d}m \cdot g \cdot \Delta h(x) = \frac{M g}{L}(2x - L)\,\mathrm{d}x
\end{equation*}


在下半段区间 $[\frac{L}{2}, L]$ 上累加积分：


\begin{equation*}
W = \int_{L/2}^L \frac{M g}{L}(2x - L)\,\mathrm{d}x = \frac{M g}{L} \left[x^2 - L x\right]_{L/2}^L
\end{equation*}



\begin{equation*}
= \frac{M g}{L} \left[(L^2 - L^2) - \left(\frac{L^2}{4} - \frac{L^2}{2}\right)\right] = \frac{M g}{L} \left(0 - \left(-\frac{L^2}{4}\right)\right) = \frac{1}{4}M g L
\end{equation*}


\textbf{[2]} 解法[2]（质心势能增量法）：

整个系统中，上半段铁链（长度 $\frac{L}{2}$，质量 $\frac{M}{2}$）始终未动，势能不变。

仅考虑下半段铁链（长度 $\frac{L}{2}$，质量 $m_2 = \frac{M}{2}$）：

拉动前，下半段铁链分布在 $[\frac{L}{2}, L]$，其质心位于中点：


\begin{equation*}
h_{\text{初}} = \frac{\frac{L}{2} + L}{2} = \frac{3}{4}L
\end{equation*}


对折完成后，下半段与上半段重叠，分布在 $[0, \frac{L}{2}]$，其质心位于新中点：


\begin{equation*}
h_{\text{末}} = \frac{0 + \frac{L}{2}}{2} = \frac{1}{4}L
\end{equation*}


因此，下半段铁链的质心被提升的高度为：


\begin{equation*}
\Delta h_C = h_{\text{初}} - h_{\text{末}} = \frac{3}{4}L - \frac{1}{4}L = \frac{1}{2}L
\end{equation*}


外力所做的功等于下半段铁链增加的重力势能：


\begin{equation*}
W = \Delta E_p = m_2 g \Delta h_C = \left(\frac{M}{2}\right) g \left(\frac{L}{2}\right) = \frac{1}{4}M g L
\end{equation*}


\textbf{[3]} 总结对比：

微元法从局部微元严格导出积分式，质心法从宏观刚体能量守恒秒杀，两者高度契合，展现了微积分与力学原理的完美融合。
\end{solbox}


\subsection{习题 12.2: 横卧圆柱形储油罐抽油做功模型}

一横卧放置的直圆柱形储油罐，底面半径为 $R$，长度为 $L$。罐内装满了密度为 $\rho$ 的油。求将罐内的油全部从油罐顶部水平切线处的开口抽出所做的功。

\textbf{\color{DeepPurple}【思路点拨】} 建立以储油罐顶面切线为基准的垂直向下坐标系，确定深度 $x$ 处水平油面的长方形截面积 $A(x) = 2L\sqrt{2Rx - x^2}$，构建提升做功微元 $\mathrm{d}W = \rho g x A(x)\,\mathrm{d}x$，利用区间对称性与三角代换完成定积分计算


\begin{solbox}[分步推导与答题规范]
\textbf{[1]} 建立坐标系与几何截面建模：

设圆柱截面的最高点为坐标原点 $O$，铅直向下为 $x$ 轴正向。

油罐截面为圆心在 $(R, 0)$、半径为 $R$ 的圆，其方程为：


\begin{equation*}
y^2 + (x - R)^2 = R^2 \implies y = \sqrt{R^2 - (x - R)^2} = \sqrt{2Rx - x^2}
\end{equation*}


油罐全深为 $x \in [0, 2R]$。在水深 $x$ 处，储油罐的水平切片为一个长方形。

长方形的长为储油罐的柱长 $L$，宽为截面圆在深度 $x$ 处的水平弦长 $2y = 2\sqrt{2Rx - x^2}$。

因此，深度 $x$ 处的水平切片截面积为：


\begin{equation*}
A(x) = 2L\sqrt{2Rx - x^2}
\end{equation*}


\textbf{[2]} 构建抽油做功微元：

取厚度为 $\mathrm{d}x$ 的水平薄油层，其微元体积为 $\mathrm{d}V = 2L\sqrt{2Rx - x^2}\,\mathrm{d}x$。

该薄油层受到的重力微元为 $\mathrm{d}F = \rho g\,\mathrm{d}V = 2\rho g L\sqrt{2Rx - x^2}\,\mathrm{d}x$。

将其抽升至油罐顶部开口（$x = 0$ 处）所需的提升位移为 $h(x) = x$。

抽油微元功为：


\begin{equation*}
\mathrm{d}W = \mathrm{d}F \cdot x = 2\rho g L x\sqrt{2Rx - x^2}\,\mathrm{d}x
\end{equation*}


\textbf{[3]} 对称性换元求解定积分：

全罐油全部抽空所做的总功为：


\begin{equation*}
W = 2\rho g L \int_0^{2R} x \sqrt{R^2 - (x - R)^2}\,\mathrm{d}x
\end{equation*}


作平移变量代换，令 $u = x - R$，则 $x = u + R$，$\mathrm{d}x = \mathrm{d}u$。积分限从 $x \in [0, 2R]$ 变为 $u \in [-R, R]$：


\begin{equation*}
\int_0^{2R} x \sqrt{R^2 - (x - R)^2}\,\mathrm{d}x = \int_{-R}^R (u + R)\sqrt{R^2 - u^2}\,\mathrm{d}u
\end{equation*}


拆分为两项：


\begin{equation*}
= \int_{-R}^R u\sqrt{R^2 - u^2}\,\mathrm{d}u + R \int_{-R}^R \sqrt{R^2 - u^2}\,\mathrm{d}u
\end{equation*}


由于第一项被积函数为对称区间上的奇函数，其积分为零；

第二项定积分 $\int_{-R}^R \sqrt{R^2 - u^2}\,\mathrm{d}u$ 为半径为 $R$ 的半圆面积，即 $\frac{1}{2}\pi R^2$。

故定积分值为：


\begin{equation*}
0 + R \cdot \frac{1}{2}\pi R^2 = \frac{1}{2}\pi R^3
\end{equation*}


代入总功公式：


\begin{equation*}
W = 2\rho g L \cdot \left(\frac{1}{2}\pi R^3\right) = \pi \rho g L R^3
\end{equation*}

\end{solbox}


\subsection{习题 12.3: 倾斜放置的矩形水闸平板静水总压力分析}

一矩形水闸板宽为 $b = 2\text{ m}$，沿水闸斜面长度为 $a = 4\text{ m}$，闸板与水平面成倾角 $\alpha = 30^\circ$ 倾斜放置。已知闸板的上边缘恰好位于水面处，求作用在闸板单侧的静水总压力（水的密度 $\rho = 1000\text{ kg/m}^3$，$g = 9.8\text{ m/s}^2$）。

\textbf{\color{DeepPurple}【思路点拨】} 沿倾斜闸板表面建立斜向一维坐标，求出对应垂直水深与面积微元，列出作用在斜微元条上的垂直压力微元，通过定积分完成求解，并应用形心深度定理进行直接验算


\begin{solbox}[分步推导与答题规范]
\textbf{[1]} 建立斜向坐标系与微元受力分析：

沿倾斜水闸板表面向下建立一维坐标轴 $s$，以闸板上边缘（水面交线）为原点 $s = 0$。

闸板沿斜面延展范围为 $s \in [0, a]$。

在斜向坐标 $s$ 处，对应的铅直水深为：


\begin{equation*}
x(s) = s \sin\alpha = s \sin 30^\circ = \frac{1}{2}s
\end{equation*}


在坐标 $s$ 处取斜向宽度为 $\mathrm{d}s$ 的微元狭条，由于闸板为矩形，水平宽度恒为 $b$。

该狭条在板面上的面积微元为：


\begin{equation*}
\mathrm{d}S = b\,\mathrm{d}s
\end{equation*}


\textbf{[2]} 构建静水压力微元与定积分：

在水深 $x(s)$ 处，液体产生的压强垂直于闸板表面，大小为：


\begin{equation*}
p(s) = \rho g x(s) = \rho g s \sin\alpha
\end{equation*}


该斜微元受到的水压力微元为：


\begin{equation*}
\mathrm{d}P = p(s)\,\mathrm{d}S = (\rho g s \sin\alpha) \cdot (b\,\mathrm{d}s) = \rho g b \sin\alpha \cdot s\,\mathrm{d}s
\end{equation*}


在整个闸板斜长 $[0, a]$ 上积分：


\begin{equation*}
P = \int_0^a \rho g b \sin\alpha \cdot s\,\mathrm{d}s = \rho g b \sin\alpha \left[\frac{s^2}{2}\right]_0^a = \frac{1}{2}\rho g b a^2 \sin\alpha
\end{equation*}


\textbf{[3]} 代入数值计算与形心法验证：

代入题设数据 $b = 2\text{ m}, a = 4\text{ m}, \sin 30^\circ = \frac{1}{2}, \rho = 1000\text{ kg/m}^3, g = 9.8\text{ m/s}^2$：


\begin{equation*}
P = \frac{1}{2} \times 1000 \times 9.8 \times 2 \times 4^2 \times \frac{1}{2} = 9800 \times 8 = 78400\text{ N} = 78.4\text{ kN}
\end{equation*}


形心法验证：

矩形闸板的几何形心位于斜边中点 $\bar{s} = \frac{a}{2} = 2\text{ m}$ 处。

形心的铅直水深为 $\bar{x} = \bar{s}\sin 30^\circ = 2 \times \frac{1}{2} = 1\text{ m}$。

闸板的总表面积为 $S = a \cdot b = 4 \times 2 = 8\text{ m}^2$。

由形心压强公式：


\begin{equation*}
P = p(\bar{x}) \cdot S = (\rho g \bar{x}) \cdot S = (1000 \times 9.8 \times 1) \times 8 = 78400\text{ N} = 78.4\text{ kN}
\end{equation*}


验算完全一致。
\end{solbox}


\subsection{习题 12.4: 铅直半圆形平板静水压力与形心定理秒杀}

一半径为 $R$ 的均质半圆形平板铅直浸入水中，其直径边恰好与水面齐平。求该半圆形平板单侧所受到的静水总压力。

\textbf{\color{DeepPurple}【思路点拨】} 解法[1]微元积分法，建立水面为横轴、向下为纵轴的坐标系，求出水深 $x$ 处的水平弦长，列出微元面积与压力微元并积分；解法[2]形心定理法，应用半圆盘的形心高度公式，由形心处静水压强乘以半圆面积直接得出结果


\begin{solbox}[分步推导与答题规范]
\textbf{[1]} 解法[1]（定积分微元法）：

取水面与半圆直径的交线为 $y$ 轴，直径中点为原点 $O$，铅直向下为 $x$ 轴正向。

半圆的边界曲线方程为 $x^2 + y^2 = R^2$ ($x \geqslant 0$)，即 $y = \sqrt{R^2 - x^2}$。

在水深 $x \in [0, R]$ 处，水平截线两端对称分布，宽度为 $w(x) = 2y = 2\sqrt{R^2 - x^2}$。

取厚度为 $\mathrm{d}x$ 的水平狭长矩形微元条，面积微元为 $\mathrm{d}S = 2\sqrt{R^2 - x^2}\,\mathrm{d}x$。

水深 $x$ 处的静水压强为 $p(x) = \rho g x$。水压力微元为：


\begin{equation*}
\mathrm{d}P = p(x)\,\mathrm{d}S = \rho g x \cdot 2\sqrt{R^2 - x^2}\,\mathrm{d}x
\end{equation*}


在区间 $[0, R]$ 上进行定积分：


\begin{equation*}
P = \int_0^R \rho g \cdot 2x\sqrt{R^2 - x^2}\,\mathrm{d}x = -\rho g \int_0^R (R^2 - x^2)^{1/2}\,\mathrm{d}(R^2 - x^2)
\end{equation*}



\begin{equation*}
= -\rho g \left[\frac{2}{3}(R^2 - x^2)^{3/2}\right]_0^R = -\rho g \left(0 - \frac{2}{3}R^3\right) = \frac{2}{3}\rho g R^3
\end{equation*}


\textbf{[2]} 解法[2]（形心定理秒杀）：

根据平面半圆盘的几何重心公式，半径为 $R$ 的半圆盘对其直径边的形心距离为：


\begin{equation*}
\bar{x} = \frac{4R}{3\pi}
\end{equation*}


半圆盘的总面积为：


\begin{equation*}
S = \frac{1}{2}\pi R^2
\end{equation*}


根据静水压力形心定理，单侧总压力等于形心深度处的液体压强与总面积的乘积：


\begin{equation*}
P = p(\bar{x}) \cdot S = (\rho g \bar{x}) \cdot S
\end{equation*}


代入表达式：


\begin{equation*}
P = \rho g \left(\frac{4R}{3\pi}\right) \cdot \left(\frac{1}{2}\pi R^2\right) = \frac{2}{3}\rho g R^3
\end{equation*}


\textbf{[3]} 分析与总结：

形心法消除了定积分凑微分过程，特别是 $\pi$ 在分母与分子中自然对消，是考研中极具威力的高速验算技巧。
\end{solbox}


\subsection{习题 12.5: 均质细棒绕端点的转动惯量与平行轴定理验证}

设有一均质细棒长为 $L$，质量为 $M$。求：(1) 该细棒关于垂直于细棒且过其端点的轴的转动惯量 $I_A$；(2) 该细棒关于垂直于细棒且过其中点的轴的转动惯量 $I_C$；(3) 验证平行轴定理 $I_A = I_C + M d^2$。

\textbf{\color{DeepPurple}【思路点拨】} 根据转动惯量微元定义 $\mathrm{d}I = r^2\,\mathrm{d}m$，分别建立以端点和中心为原点的直角坐标系，将质量微元表示为线密度乘以位移微分，实施定积分推导，最后将两转动惯量与质心平移距离代入平行轴定理公式验证一致性


\begin{solbox}[分步推导与答题规范]
\textbf{[1]} 求细棒绕端点的转动惯量 $I_A$：

取细棒一端 $A$ 为原点 $O$，沿细棒方向为 $x$ 轴。细棒分布在 $x \in [0, L]$。

细棒线密度为 $\mu = \frac{M}{L}$。

在坐标 $x$ 处取长度为 $\mathrm{d}x$ 的质量微元 $\mathrm{d}m = \mu\,\mathrm{d}x = \frac{M}{L}\,\mathrm{d}x$。

该微元到转轴 $A$ 的旋转距离为 $r = x$。

转动惯量微元为：


\begin{equation*}
\mathrm{d}I_A = x^2\,\mathrm{d}m = \frac{M}{L} x^2\,\mathrm{d}x
\end{equation*}


定积分计算：


\begin{equation*}
I_A = \int_0^L \frac{M}{L} x^2\,\mathrm{d}x = \frac{M}{L} \left[\frac{x^3}{3}\right]_0^L = \frac{1}{3}M L^2
\end{equation*}


\textbf{[2]} 求细棒绕中点的转动惯量 $I_C$：

取细棒中心 $C$（质心）为原点，沿细棒方向为 $x$ 轴。细棒分布在 $x \in [-\frac{L}{2}, \frac{L}{2}]$。

在坐标 $x$ 处取质量微元 $\mathrm{d}m = \frac{M}{L}\,\mathrm{d}x$。

该微元到质心转轴 $C$ 的旋转距离为 $r = |x|$。

转动惯量微元为 $\mathrm{d}I_C = x^2\,\mathrm{d}m = \frac{M}{L} x^2\,\mathrm{d}x$。

定积分计算：


\begin{equation*}
I_C = \int_{-L/2}^{L/2} \frac{M}{L} x^2\,\mathrm{d}x = 2 \frac{M}{L} \left[\frac{x^3}{3}\right]_0^{L/2} = \frac{2M}{3L} \left(\frac{L^3}{8}\right) = \frac{1}{12}M L^2
\end{equation*}


\textbf{[3]} 平行轴定理验证：

平行轴定理表明：刚体绕任意轴的转动惯量，等于刚体绕通过其质心且与该轴平行的轴的转动惯量，加上刚体质量与两轴间垂直距离平方的乘积：


\begin{equation*}
I = I_C + M d^2
\end{equation*}


在本题中，端点轴 $A$ 与质心轴 $C$ 相互平行，两轴间的垂直几何距离为 $d = \frac{L}{2}$。

代入公式右边：


\begin{equation*}
I_C + M d^2 = \frac{1}{12}M L^2 + M \left(\frac{L}{2}\right)^2 = \frac{1}{12}M L^2 + \frac{1}{4}M L^2 = \left(\frac{1}{12} + \frac{3}{12}\right)M L^2 = \frac{4}{12}M L^2 = \frac{1}{3}M L^2
\end{equation*}


恰好等于 $I_A$，平行轴定理严格成立。
\end{solbox}


\subsection{习题 12.6: 非线性变弹性力做功与势能分析}

某非线性减震弹性元件的恢复力与位移 $x$ 满足关系 $F(x) = k_1 x + k_2 x^3$（其中 $k_1 > 0, k_2 > 0$）。求将该元件自平衡位置 $x = 0$ 缓慢拉伸至位移 $x = X$ 过程中外力克服弹性力所做的功，并在 $X = 0.2\text{ m}, k_1 = 200\text{ N/m}, k_2 = 5000\text{ N/m}^3$ 条件下计算具体功的数值。

\textbf{\color{DeepPurple}【思路点拨】} 根据一元积分变力沿直线做功的微元模型，在位移 $[x, x+\mathrm{d}x]$ 上外力等于克服的弹性力大小，写出功微元 $\mathrm{d}W = (k_1 x + k_2 x^3)\,\mathrm{d}x$，在 $[0, X]$ 上定积分求解，随后代入工程参数计算实际做功数值


\begin{solbox}[分步推导与答题规范]
\textbf{[1]} 建立变力做功的定积分微元：

当弹性元件被拉长 $x$ 时，内弹力大小为 $F(x) = k_1 x + k_2 x^3$。

缓慢拉伸过程中，外力大小与弹性恢复力大小处处平衡，故外力函数为：


\begin{equation*}
F_{\text{外}}(x) = k_1 x + k_2 x^3
\end{equation*}


在微小位移区间 $[x, x + \mathrm{d}x]$ 上，外力所做的功微元为：


\begin{equation*}
\mathrm{d}W = F_{\text{外}}(x)\,\mathrm{d}x = (k_1 x + k_2 x^3)\,\mathrm{d}x
\end{equation*}


\textbf{[2]} 定积分计算功的代数表达式：

将元件从 $x = 0$ 拉伸到 $x = X$ 所做的总功为：


\begin{equation*}
W = \int_0^X (k_1 x + k_2 x^3)\,\mathrm{d}x = \left[\frac{1}{2}k_1 x^2 + \frac{1}{4}k_2 x^4\right]_0^X = \frac{1}{2}k_1 X^2 + \frac{1}{4}k_2 X^4
\end{equation*}


\textbf{[3]} 代入工程数值计算：

将参数 $X = 0.2\text{ m}, k_1 = 200\text{ N/m}, k_2 = 5000\text{ N/m}^3$ 代入：


\begin{equation*}
W = \frac{1}{2} \times 200 \times (0.2)^2 + \frac{1}{4} \times 5000 \times (0.2)^4
\end{equation*}



\begin{equation*}
= 100 \times 0.04 + 1250 \times 0.0016 = 4 + 2 = 6\text{ J}
\end{equation*}


外力克服弹性恢复力共做功 $6\text{ 焦耳}$。
\end{solbox}


\subsection{习题 12.7: 线性供求函数下垄断企业利润最大化与税收效应}

已知某垄断企业的边际收益函数为 $MR(Q) = 140 - 4Q$，边际成本函数为 $MC(Q) = 20 + 2Q$（万元/千件），固定成本为 $C_0 = 100$ 万元。求：(1) 企业实现最大总利润时的产量 $Q^*$ 及此时的最大总利润；(2) 若政府对每千件产品征收税额 $t = 12$ 万元的从量税，求税后企业的最优产量与最大利润。

\textbf{\color{DeepPurple}【思路点拨】} 根据边际收益等于边际成本确定利润最大化的产量；利用定积分分别还原总收益函数与总成本函数，求差获得总利润；征税后边际成本增加 $t$，重新求解均衡产量与税后最大利润


\begin{solbox}[分步推导与答题规范]
\textbf{[1]} 无税条件下的利润最大化产量与利润：

总利润取得最大值的必要条件是边际收益等于边际成本：


\begin{equation*}
MR(Q) = MC(Q) \implies 140 - 4Q = 20 + 2Q
\end{equation*}


解得最优产量为：


\begin{equation*}
6Q = 120 \implies Q^* = 20\text{ 千件}
\end{equation*}


由微积分还原总收益与总成本函数：


\begin{equation*}
R(Q) = \int_0^Q MR(t)\,\mathrm{d}t = \int_0^Q (140 - 4t)\,\mathrm{d}t = 140Q - 2Q^2
\end{equation*}



\begin{equation*}
C(Q) = C_0 + \int_0^Q MC(t)\,\mathrm{d}t = 100 + \int_0^Q (20 + 2t)\,\mathrm{d}t = 100 + 20Q + Q^2
\end{equation*}


总利润函数为：


\begin{equation*}
L(Q) = R(Q) - C(Q) = (140Q - 2Q^2) - (100 + 20Q + Q^2) = -3Q^2 + 120Q - 100
\end{equation*}


代入 $Q^* = 20$：


\begin{equation*}
L(20) = -3 \times 20^2 + 120 \times 20 - 100 = -1200 + 2400 - 100 = 1100\text{ 万元}
\end{equation*}


\textbf{[2]} 征收从量税后的边际成本与新最优产量：

政府对每千件产品征税 $t = 12$ 万元后，企业的实际边际成本增加为：


\begin{equation*}
MC_t(Q) = MC(Q) + t = 20 + 2Q + 12 = 32 + 2Q
\end{equation*}


新的利润最大化条件为：


\begin{equation*}
MR(Q) = MC_t(Q) \implies 140 - 4Q = 32 + 2Q
\end{equation*}


解得税后最优产量为：


\begin{equation*}
6Q = 108 \implies Q_t^* = 18\text{ 千件}
\end{equation*}


\textbf{[3]} 税后最大总利润计算：

税后总成本函数为 $C_t(Q) = C(Q) + tQ = 100 + 32Q + Q^2$。

税后总利润函数为：


\begin{equation*}
L_t(Q) = R(Q) - C_t(Q) = (140Q - 2Q^2) - (100 + 32Q + Q^2) = -3Q^2 + 108Q - 100
\end{equation*}


代入 $Q_t^* = 18$：


\begin{equation*}
L_t(18) = -3 \times 18^2 + 108 \times 18 - 100 = -972 + 1944 - 100 = 872\text{ 万元}
\end{equation*}


企业总利润因税收减少了 $1100 - 872 = 228\text{ 万元}$。
\end{solbox}


\subsection{习题 12.8: 非线性需求函数下的消费者剩余定积分计算}

某高端消费品市场的需求函数为 $P = \frac{120}{\sqrt{2Q+1}}$（万元/单位）。当市场定价为 $P_0 = 24$ 万元时，求：(1) 消费者的市场购买量 $Q_0$；(2) 市场交易中的消费者剩余 $CS$。

\textbf{\color{DeepPurple}【思路点拨】} 由市场销售价格等于需求价格求出均衡消费量 $Q_0$；根据消费者剩余的定积分定义式 $CS = \int_0^{Q_0} [P_D(Q) - P_0]\,\mathrm{d}Q$，通过凑微分计算含无理根式的定积分


\begin{solbox}[分步推导与答题规范]
\textbf{[1]} 求解市场均衡消费量 $Q_0$：

当市场单价为 $P_0 = 24$ 时，代入需求函数：


\begin{equation*}
\frac{120}{\sqrt{2Q_0 + 1}} = 24
\end{equation*}


方程两边除以 $24$：


\begin{equation*}
\sqrt{2Q_0 + 1} = 5 \implies 2Q_0 + 1 = 25 \implies Q_0 = 12\text{ 单位}
\end{equation*}


\textbf{[2]} 列出消费者剩余定积分表达式：

消费者剩余是需求价格线在 $[0, Q_0]$ 上超出实际支付价格 $P_0$ 的累积积分：


\begin{equation*}
CS = \int_0^{Q_0} [P(Q) - P_0]\,\mathrm{d}Q = \int_0^{12} \left(\frac{120}{\sqrt{2Q + 1}} - 24\right)\,\mathrm{d}Q
\end{equation*}



\begin{equation*}
= \int_0^{12} \frac{120}{\sqrt{2Q + 1}}\,\mathrm{d}Q - \int_0^{12} 24\,\mathrm{d}Q
\end{equation*}


\textbf{[3]} 定积分计算求值：

计算第一项定积分，凑微分：


\begin{equation*}
\int_0^{12} \frac{120}{\sqrt{2Q + 1}}\,\mathrm{d}Q = 60 \int_0^{12} (2Q + 1)^{-1/2}\,\mathrm{d}(2Q + 1)
\end{equation*}



\begin{equation*}
= 60 \left[2(2Q + 1)^{1/2}\right]_0^{12} = 120 (\sqrt{25} - \sqrt{1}) = 120 (5 - 1) = 480
\end{equation*}


计算第二项矩形支出面积：


\begin{equation*}
P_0 Q_0 = 24 \times 12 = 288
\end{equation*}


相减即得消费者剩余：


\begin{equation*}
CS = 480 - 288 = 192\text{ 万元}
\end{equation*}

\end{solbox}

\end{document}
